\subsection{声强与声压、质元速度怎样联系}
声音传播时，空气质元在平衡位置附近往复运动，并不是从声源一路飞到听者。取沿 $+x$ 传播的平面简谐声波，位移为
\[
 \xi(x,t)=\xi_m\cos(kx-\omega t),\qquad \omega=ck.
\]
质元速度是对时间求偏导，而波速是波形的传播速度，二者不能混淆：
\[
 v_{\rm part}=\partial_t\xi=\omega\xi_m\sin(kx-\omega t).
\]
小形变的体压缩量为 $-\partial_x\xi$；用绝热体积模量 $B_s=\rho_0c^2$，声压扰动为
\[
 p'=-B_s\partial_x\xi
 =\rho_0c\omega\xi_m\sin(kx-\omega t)
 =\rho_0c\,v_{\rm part}.
\]
这里 $p'$ 是总压强相对平衡值的微小变化，可以为正或负；$\rho_0$ 为平衡密度。对于这列右行波，声压与质元速度同相。截面上每单位面积传递的瞬时功率为 $p'v_{\rm part}$，取一个周期的平均，得到
\[
 I=\langle p'v_{\rm part}\rangle
 =\frac12\rho_0c\omega^2\xi_m^2
 =\frac{p_m^2}{2\rho_0c}
 =\frac{p_{\rm rms}^2}{\rho_0c}.
\]
因正弦平方的周期平均为 $1/2$，有效声压 $p_{\rm rms}=p_m/\sqrt2$。所以声强与振幅的平方成正比，而不是与振幅本身成正比。上述声压与声强关系适用于单向平面行波；在驻波中声压和速度的时空关系不同，不能照搬。

\subsection{先算强度，再换成分贝}
球面面积随距离平方增加。在无吸收且各向同性辐射的条件下，每个球面通过的总功率都为 $P$，所以 $4\pi r^2I=P$。这不是能量丢失，而是同一功率分散到更大的面积。距离从 $r_1$ 变成 $r_2$ 时，
\[
 L_I(r_2)-L_I(r_1)
 =10\log_{10}\frac{I(r_2)}{I(r_1)}
 =-20\log_{10}\frac{r_2}{r_1}.
\]
距离加倍，声强减为四分之一，声强级下降约 $6.02\,\mathrm{dB}$。这与“两个独立相同声源使声强级增加约 $3.01\,\mathrm{dB}$”并不矛盾，前者改变传播面积，后者改变总入射强度。

\begin{origexamplebox}{不同距离、不同声强级的两个独立声源}\label{ch12:ex:teaching-sound-distance}
两个可视为点声源的机器各向同性发声。在距机器甲 $2\,\mathrm m$ 处，它单独产生 $80\,\mathrm{dB}$ 的声强级；机器乙在观察点单独产生 $70\,\mathrm{dB}$。观察点距甲 $4\,\mathrm m$。忽略反射和吸收，二者互不相干。求同时工作时观察点的总声强级。
\end{origexamplebox}
\begin{xesolution}
先把甲的测量位置换到观察点。距离加倍，强度减为四分之一，故
\[
 I_A=\frac14I_0 10^{80/10}=2.5\times10^{-5}\,\mathrm{W/m^2}.
\]
乙在该点的强度为 $I_B=I_0 10^{70/10}=10^{-5}\,\mathrm{W/m^2}$。互不相干时，平均交叉项为零，才有 $I=I_A+I_B$，所以
\[
 L=10\log_{10}(3.5\times10^7)\approx75.44\,\mathrm{dB}.
\]
结果略高于甲单独在该点的 $73.98\,\mathrm{dB}$，不是两个分贝数的算术和。若乙关闭，保留 $I_A$ 即回到甲单独的结果。
\end{xesolution}

若两列同频同向平面声波保持固定相位差 $\delta$，应先叠加声压，再平方取平均，得到
\[
 I=I_1+I_2+2\sqrt{I_1I_2}\cos\delta.
\]
等强同相时局部强度为 $4I_1$，反相时可以相消；若相位不断随机变化，平均 $\cos\delta=0$，才恢复非相干相加。驻波、不同传播方向和复杂声场还应计算矢量能流，不能把这条同向行波公式当作任意声场的通式。
